\chapter{Kết luận và Hướng phát triển}

\section{Các vấn đề thách thức và Giải pháp}

Trong quá trình thiết kế và phát triển engine cờ tướng Lingine, nhóm nghiên cứu đã đối mặt và giải quyết thành công một số thách thức kỹ thuật lớn:

\begin{enumerate}
    \item \textbf{Hiệu ứng chân trời (Horizon Effect):}
    \begin{itemize}
        \item \textit{Thách thức:} Xảy ra khi engine phát hiện một đòn chiến thuật bất lợi (ví dụ bị bắt Xe) ở giới hạn độ sâu tìm kiếm. Để trì hoãn việc này ra sau ``chân trời'' tìm kiếm, engine có thể thực hiện các nước đi thí quân vô nghĩa nhằm kéo dài độ sâu, dẫn đến đánh giá sai lầm nghiêm trọng.
        \textit{Giải pháp:} Triển khai bộ tìm kiếm tĩnh Quiescence Search kết hợp với kiểm tra chiếu tướng mở rộng (Check Extensions). Quiescence Search tiếp tục duyệt tất cả các nước đi ăn quân cho đến khi thế trận hoàn toàn bình ổn tĩnh, đảm bảo điểm số trả về phản ánh chính xác thực tế thế cờ.
    \end{itemize}
    \item \textbf{Tranh chấp đồng bộ hóa trong giao thức UCI:}
    \begin{itemize}
        \item \textit{Thách thức:} Lệnh ngắt \texttt{stop} từ người dùng có thể gửi đến bất kỳ lúc nào qua luồng vào chuẩn khi engine đang tính toán sâu. Nếu sử dụng các cơ chế đồng bộ luồng chặn hoặc khóa Mutex truyền thống sẽ gây trễ hoặc làm treo toàn bộ ứng dụng.
        \item \textit{Giải pháp:} Áp dụng biến nguyên tử \texttt{AtomicBool} chia sẻ giữa luồng đọc lệnh và luồng tìm kiếm. Thao tác đọc ghi bất đồng bộ không chặn giúp dừng tìm kiếm tức thời trong thời gian dưới 1 mili-giây mà không gây ra hiện tượng nghẽn luồng.
    \end{itemize}
\end{enumerate}

\section{Công nghệ và Tài nguyên sử dụng}

Dự án Lingine được xây dựng và kiểm thử dựa trên các công nghệ và tài nguyên mã nguồn mở hiện đại:
\begin{itemize}
    \item \textbf{Rust Toolchain:} Trình biên dịch \texttt{rustc} và công cụ quản lý dự án \texttt{cargo}. Chúng tôi tận dụng tối đa các tối ưu hóa biên dịch của Cargo thông qua cấu hình biên dịch \texttt{opt-level = 3} và kích hoạt các chỉ thị tập lệnh CPU vector hiện đại (\texttt{-C target-cpu=native}).
    \item \textbf{Thư viện \texttt{arrayvec}:} Cho phép lưu trữ danh sách nước đi (\texttt{MoveList}) trong bộ nhớ Stack cố định thay vì bộ nhớ Heap động, tránh overhead cấp phát bộ nhớ khi sinh hàng triệu nước đi mỗi giây.
    \item \textbf{Python Suite:} Các kịch bản Python tự động thiết lập môi trường (\texttt{setup\_}\allowbreak\texttt{tools.py}) và tổ chức đấu giải ước lượng Elo (\texttt{run\_}\allowbreak\texttt{match.py}, \texttt{run\_}\allowbreak\texttt{gauntlet.py}, \texttt{run\_}\allowbreak\texttt{historical\_}\allowbreak\texttt{evals.py}).
    \item \textbf{Trình quản lý Sylvan:} Công cụ \textit{Sylvan} được sử dụng làm giao diện người dùng đồ họa (GUI) để kiểm thử thủ công và \textit{Sylvan-CLI} đóng vai trò làm bộ điều phối giải đấu thông qua giao thức UCI chuẩn mực.
    \item \textbf{Các engine cờ tướng đối chiếu:} \textit{Fairy-Stockfish} và \textit{Pikafish} làm các đối thủ có giới hạn sức mạnh chuẩn để chạy giải đấu ước lượng ELO và là công cụ kiểm thử tính đúng đắn của luật chơi.
    \item \textbf{Dự án cờ vua mã nguồn mở bằng Rust:} Tham khảo cấu trúc kiến trúc mã nguồn tối ưu hiệu năng từ hai dự án \textit{Reckless} và \textit{Viridithas} viết bằng Rust.
\end{itemize}

\section{Kết luận}

Dự án nghiên cứu thiết kế và triển khai engine cờ tướng Lingine đã hoàn thành đầy đủ tất cả các mục tiêu đề ra ban đầu:
\begin{itemize}
    \item Xây dựng thành công một engine cờ tướng hoàn chỉnh bằng ngôn ngữ Rust, đạt độ chính xác luật chơi tuyệt đối được xác minh qua các bộ kiểm thử PERFT chuẩn quốc tế.
    \item Triển khai cấu trúc dữ liệu Bitboard 128-bit kết hợp với Magic Bitboards giúp tối ưu hóa cực đoan tốc độ sinh nước đi, đạt hiệu suất xử lý hàng triệu nút trạng thái trên giây.
    \item Tích hợp hàm đánh giá tĩnh tapered score hiện đại, tự động tối ưu hóa các tham số thông qua thuật toán Texel Tuning.
    \item Triển khai bộ thuật toán tìm kiếm nâng cao Negamax tích hợp PVS, bảng băm TT và các kỹ thuật tỉa nhánh hiện đại giúp tăng độ sâu chiến thuật và đạt mức Elo thi đấu ước lượng tối đa khoảng 2244.
\end{itemize}

\section{Hướng phát triển tương lai}

Để tiếp tục nâng cao sức mạnh và tính thực tiễn của engine, nhóm phát triển đề xuất ba hướng nghiên cứu trọng tâm tiếp theo:

\subsection{Tăng tính đa dạng trong lối chơi bằng Softmax Sampling}

Trong phiên bản hiện tại, cơ chế ra quyết định của Lingine mang tính xác định tuyệt đối (deterministic) bằng cách luôn ưu tiên chọn nước đi có điểm số đánh giá cao nhất (\texttt{bestmove}). Mặc dù tối ưu về mặt lý thuyết, cách tiếp cận này khiến lối chơi của engine trở nên rập khuôn và rất dễ bị đối thủ bắt bài, đặc biệt là trong các ván đấu tập hoặc ở giai đoạn khai cuộc.

Để khắc phục, Lingine sẽ tích hợp kỹ thuật lấy mẫu Softmax (Softmax Sampling) nhằm xác suất hóa các lựa chọn nước đi. Thay vì chỉ chọn giá trị lớn nhất (Argmax), thuật toán sẽ chuyển đổi điểm số $Q(s, a)$ của mọi nước đi hợp lệ thành một phân phối xác suất:

\begin{equation}
P(a_i | s) = \frac{e^{Q(s, a_i) / \tau}}{\sum_j e^{Q(s, a_j) / \tau}}
\end{equation}

Trong đó, $\tau$ là tham số nhiệt độ (Temperature) đóng vai trò kiểm soát độ nhiễu. Khi $\tau \to 0$, hàm hội tụ về Argmax (chỉ đánh nước tốt nhất). Khi $\tau$ tăng dần, các nước đi yếu hơn cũng có xác suất được chọn, tạo ra lối chơi sáng tạo và khó đoán.

\begin{figure}[htbp]
    \centering
    \begin{tikzpicture}
        % Axes
        \draw[thick, ->] (0,0) -- (8,0) node[right] {Nước đi (Chất lượng giảm dần)};
        \draw[thick, ->] (0,0) -- (0,5) node[above] {Xác suất chọn $P(a_i)$};
        
        % Argmax (tau -> 0)
        \draw[blue, ultra thick] (1, 4.5) -- (1, 0);
        \fill[blue] (1, 4.5) circle (2.5pt);
        \node[blue, right] at (1.1, 4.5) {$\tau \to 0$ (Argmax)};
        
        % Softmax moderate (tau = 0.5)
        \draw[red, thick, smooth] plot coordinates {(0.5, 0.5) (1, 3.5) (2, 1.5) (3, 0.8) (4, 0.4) (6, 0.1)};
        \node[red, right] at (2.1, 1.8) {$\tau = 0.5$ (Cân bằng)};
        
        % Softmax high (tau -> inf)
        \draw[green!60!black, thick, dashed, smooth] plot coordinates {(0.5, 1.2) (1, 1.1) (3, 0.9) (5, 0.8) (7, 0.7)};
        \node[green!60!black, right] at (5.1, 1.0) {$\tau \to \infty$ (Ngẫu nhiên)};
        
        % Labels on X-axis
        \draw (1, 0.1) -- (1, -0.1) node[below] {Tốt nhất};
        \draw (3, 0.1) -- (3, -0.1) node[below] {Tốt nhì};
        \draw (5, 0.1) -- (5, -0.1) node[below] {Bình thường};
    \end{tikzpicture}
    \caption{Phân phối xác suất chọn nước đi thay đổi theo tham số nhiệt độ $\tau$}
    \label{fig:softmax_sampling}
\end{figure}


\subsection{Tích hợp Mạng nơ-ron NNUE (Efficiently Updatable Neural Networks)}

\begin{figure}[htbp]
    \centering
    \begin{tikzpicture}[
        node distance=1.5cm and 2cm,
        neuron/.style={circle, draw, minimum size=0.8cm, inner sep=0pt},
        active/.style={fill=red!30},
        hidden/.style={fill=blue!20},
        output/.style={fill=green!30}
    ]
        % Input Layer (Sparse)
        \node[neuron] (i1) at (0, 3) {};
        \node[neuron, active] (i2) at (0, 2) {};
        \node[neuron] (i3) at (0, 1) {};
        \node at (0, 0.3) {$\vdots$};
        \node[neuron, active] (in) at (0, -1) {};
        
        \node[above=0.5cm of i1, text width=2.5cm, align=center] {Lớp đầu vào\\(Đặc trưng thưa)};
        \node[left=0.2cm of i2] {Quân di chuyển $\to$};
        
        % Accumulator (Hidden Layer 1)
        \node[neuron, hidden] (a1) at (4, 2.5) {};
        \node[neuron, hidden] (a2) at (4, 1.5) {};
        \node[neuron, hidden] (a3) at (4, 0.5) {};
        \node at (4, -0.2) {$\vdots$};
        \node[neuron, hidden] (am) at (4, -1) {};
        
        \node[above=1cm of a1, text width=3cm, align=center] {Bộ tích lũy\\(Cập nhật vi phân)};
        
        % Output Layer
        \node[neuron, output] (out) at (8, 0.75) {};
        \node[above=0.5cm of out, text width=2cm, align=center] {Điểm số đánh giá};
        
        % Connections Input -> Accumulator (Highlighting incremental)
        \foreach \i in {i2, in} {
            \foreach \a in {a1, a2, a3, am} {
                \draw[->, thick, red] (\i) -- (\a);
            }
        }
        \draw[->, gray, dashed] (i1) -- (a1);
        \draw[->, gray, dashed] (i3) -- (a2);
        
        % Connections Accumulator -> Output
        \foreach \a in {a1, a2, a3, am} {
            \draw[->, thick, blue] (\a) -- (out);
        }
    \end{tikzpicture}
    \caption{Cơ chế cập nhật vi phân của mạng NNUE}
    \label{fig:nnue_architecture}
\end{figure}

Hàm đánh giá tĩnh của Lingine hiện đang dựa trên mô hình tuyến tính (Material + PST). Dù đem lại ưu thế về tốc độ xử lý, cấu trúc tuyến tính này không thể nắm bắt được những tương tác phi tuyến tính phức tạp trong cờ tướng (ví dụ: Xe-Pháo-Mã phối hợp).

Thay thế hàm tĩnh bằng mạng nơ-ron NNUE dạng nông. Điểm làm nên sức mạnh của NNUE là \textit{kiến trúc lớp đầu vào cập nhật vi phân}. Thay vì tính toán lại toàn bộ bàn cờ từ đầu sau mỗi nước đi, hệ thống chỉ nhân ma trận cho các quân cờ vừa dịch chuyển (do\_move/undo\_move), giúp duy trì tốc độ duyệt (NPS) cực cao.


\subsection{Tối ưu hóa tham số bằng thuật toán SPSA}

Hàng trăm tham số tìm kiếm (Search Parameters) như hệ số cắt tỉa LMR, Futility margins vẫn đang phải tinh chỉnh thủ công. Không thể dùng Gradient Descent truyền thống vì đây là hệ thống hộp đen (không tính được đạo hàm).

Áp dụng thuật toán \textit{Simultaneous Perturbation Stochastic Approximation (SPSA)}. Thuật toán sẽ tác động nhiễu loạn ngẫu nhiên lên toàn bộ vector tham số $\theta$ cùng lúc, tạo ra hai bộ tham số $\theta + \Delta$ và $\theta - \Delta$. Bằng cách cho hai phiên bản engine này tự đấu với nhau, hệ thống ước lượng được gradient và tự động cập nhật để tìm ra mốc cắt tỉa tối ưu nhất.

\begin{figure}[htbp]
    \centering
    \begin{tikzpicture}[
        node distance=1.5cm and 2cm,
        box/.style={rectangle, draw, rounded corners, minimum height=1.2cm, text centered, text width=3.8cm},
        box_sq/.style={rectangle, draw, minimum height=1.2cm, text centered, text width=3.5cm},
        arrow/.style={thick, ->, >=stealth}
    ]
        % Các node (hộp)
        \node[box, fill=green!20, text width=4cm] (theta) at (0, 4) {Tham số hiện tại $\theta$};
        
        \node[box, fill=blue!10] (plus) at (-2.5, 2) {$\theta + c_k\Delta_k$};
        \node[box, fill=blue!10] (minus) at (2.5, 2) {$\theta - c_k\Delta_k$};
        
        \node[box_sq, fill=orange!20] (game) at (0, 0) {Tự đấu\\(Engine Match)};
        
        \node[box, fill=pink!20, text width=4.5cm] (grad) at (0, -2) {Ước lượng\\Gradient $\hat{g}_k(\theta)$};
        
        % Các mũi tên đường đi xuống
        \draw[arrow] (theta.south) -- (plus.north) node[midway, left, yshift=0.15cm] {+ Nhiễu $\Delta_k$};
        \draw[arrow] (theta.south) -- (minus.north) node[midway, right, yshift=0.15cm] {- Nhiễu $\Delta_k$};
        
        \draw[arrow] (plus.south) -- (game.north west);
        \draw[arrow] (minus.south) -- (game.north east);
        
        \draw[arrow] (game.south) -- (grad.north) node[midway, right=0.1cm] {Thắng/Thua/Hòa};
        
        % Mũi tên vòng lặp (Feedback loop)
        \draw[arrow, rounded corners] (grad.west) -- (-4.8, -2) -- (-4.8, 4) -- (theta.west);
        
        % Label cho vòng lặp cập nhật
        \node[anchor=south west, align=left] at (-5.0, 4.1) {Cập nhật $\theta_{k+1}$\\$\theta_k - a_k \hat{g}_k(\theta)$};
        
    \end{tikzpicture}
    \caption{Chu trình tối ưu hóa tham số hộp đen bằng thuật toán SPSA}
    \label{fig:spsa_flowchart}
\end{figure}
